\section{The original compatibility conditions at a selected member}

% BEGIN M84T1
\paragraph{M84T1. The period reading already present at an original member.}
Fix the local correspondence $h_L$ in the source and an actual
member in $\mathrm{MemberAtLayerPoint}(h_L)$. Write $D=w.D(Z)$.
For every retained coefficient index $i$ and $v\mid\ell_i$,
the inertia polynomial of $D.\mathrm{wdOf}_i(r,v)$ for a continuous
integral representation $r$, potentially crystalline at $v$, is the
inertia characteristic polynomial of its potentially semistable
Weil--Deligne realization, using the very period bridge $D.\mathrm{nbr}$
in the test. Here and below the coefficients are mapped by the selected
embedding $\iota_i:\overline{\mathbb Q}_{\ell_i}\simeq\mathbb C$.

To identify the input on this member, unpack its $c.a.\mathrm{nt}$
realization, its equality of data and its character-tuple equality.
The three projections used here are the Weil--Deligne reading
$A.\mathrm{wdT}$, the global-to-local comparison $A.\mathrm{wd}$,
and the equality $\mathrm{Au}.\mathrm{wdOf}$ at the added localization.
They are the three fields copied in the construction of a
$\mathrm{CharacterChoiceData}$ from this realization. They can be
read before constructing its other fields. The last equality identifies
the native parameter with the one read by $A.\mathrm{wd}$.
The inertia clause of $A.\mathrm{wdT}$ identifies its inertia
polynomial. The above-$\ell$ clause of $A.\mathrm{wd}$ reads it
on a bridge inducing the specified place above $v$.

Here is the comparison with the particular bridge in the test.
Any two embeddings of the chosen global algebraic closure into
$\overline{\mathbb Q}_\ell$ inducing that place differ, on their
images, by a local Galois conjugacy after identifying the two
base completions over the dense number field. Indeed the resulting
isometry of the dense algebraic subfields extends to the completions;
over the common completed base the two embeddings of its algebraic
closure are conjugate, by extension of embeddings. Transport the
Galois action on period tensors by this conjugacy, acting on the
period factor, with the conjugate restriction of $r$ on the vector
factor. This identifies the invariant modules: the subgroup fixing
a vector is carried to its conjugate. It commutes with period
Frobenius and monodromy and conjugates the two actions of each
original inertia element. In particular potential semistability
and the characteristic polynomials agree. To compare finite
coefficient models, extend both to their compositum. The matrices
of the common representation then agree entry by entry under the
fixed embeddings in $\overline{\mathbb Q}_\ell$. Extension of
coefficients commutes with taking period invariants here: a basis
for the finite coefficient extension gives a finite direct sum of
the original invariant equations. The induced scalarized period
spaces and their inertia actions thus give the same characteristic
polynomials. This argument preserves the complex embedding and
the specified inertia element while changing the auxiliary
period realization.
% END M84T1

% BEGIN M84T2
\paragraph{M84T2. The inertia polynomial kills every specified crystalline period.}
On every original member of M84T1 the full third condition in
$\mathrm{Tie4}$, namely $\mathrm{TieInePoly}$, holds. Fix the
index, place, normal finite-index local extension in the test,
continuous rank-$m$ representation, inertia element $s$, and
polynomial $q\in O_i[X]$ whose coefficient image is the native
inertia polynomial. Let $d$ be any of the tested potentially
crystalline periods. The assertion is
\[
  \sum_{j=0}^{\deg q} \mathrm{toLL}(q_j)\,
                      \mathrm{act}_{V,B}(e(s^j))d=0 .
\]
Here $e$ is the original bridge map on the decomposition group.
Potentially crystalline periods are potentially semistable under
$B_{\rm cris}\hookrightarrow B_{\rm st}$. After a finite coefficient
extension the original model and the one giving the Weil--Deligne
reading agree by M84T1. Write $A_s$ for the action of $s$ on its
$D_{\rm pst}$. Inertia fixes the maximal unramified scalar field,
so $A_s$ is linear there. It also commutes with period Frobenius
$\varphi$. In degree zero of the Weil homomorphism the Frobenius
correction and the unipotent factor used for Frobenius
semisimplification both act as the identity. Therefore the raw
inertia action on the complex period space has characteristic
polynomial $\iota_i(q)$, by M84T1. Cayley--Hamilton gives
\[
             \mathrm{wdMk}\big(q(A_s)y\big)=0
                    \qquad (y\in D_{\rm pst}).
\]

For clarity this last map evaluates a single coefficient component.
Over $\mathbb Q_\ell^{\rm ur}$ the scalar algebra of a finite
coefficient field $E'$ is $E'\otimes_{\mathbb Q_\ell}
\mathbb Q_\ell^{\rm ur}$. Its field factors are indexed by the
embeddings of the maximal unramified subfield of $E'$: the remaining
totally ramified extension stays a field over the maximal unramified
extension, by its Eisenstein uniformizer polynomial. Scalarization
is faithful on its chosen factor since it extends a field embedding.
Frobenius permutes the factors transitively and is bijective on
$D_{\rm pst}$. As in OFL7, applying the last display to every
$\varphi^j y$ detects every component of $q(A_s)y$, since
$\varphi$ commutes with $A_s$ and fixes $E'$ acting as coefficient
scalars. Hence $q(A_s)=0$ on the whole period module.

The coefficient inclusion on $V$ is injective, and remains so
after tensoring over $\mathbb Q_\ell$ with the period ring.
Likewise the inclusion of crystalline in semistable periods is
injective, since the vector space is flat. A period fixed over
the given finite extension maps into $D_{\rm pst}$, and both
inclusions intertwine $s$ and the coefficient scalars. Pulling the
polynomial identity back gives $q(\mathrm{act}(e(s)))d=0$.
On pure tensors $\mathrm{act}(e(s))^j=\mathrm{act}(e(s^j))$
by the group laws on the two factors. Expanding $q$ proves the
display with the native coefficients and every extension, rank,
matrix representation, polynomial and period in its quantifiers.
% END M84T2

% BEGIN M84T3
\paragraph{M84T3. The central action and the full Hecke eigenspace.}
The following constructs the integral-action and eigenspace parts of
the original member's sought $\mathrm{Pend}$ witness from its
localized type sections and integral class-sum test. At index $i$
put $M=D.\mathrm{Mt}_i$, $V=K\otimes_O M$, where $O=O_i$ and
$K=\operatorname{Frac}(O)$. Let $C$ be the complex space of forms,
$C_x$ the sum of the type spaces of representatives of $x$, and
$C_\emptyset$ the corresponding sum for representatives with no
Hecke point. Use the original independent type spaces and their
decomposition in the row. Suppose $\Phi:V\to C$ is a localized
type section: it is injective and semilinear for the given
embedding, $\Phi(V_x)\subset C_x$, and the difference between a
rational form and $\Phi$ of its projection is in $C_\emptyset$.
Assume the finite group action preserves each type space and
integral central elements send each distinguished rational lattice
generator into the lattice of rational forms. These are precisely
the section, type stability, and $\mathrm{CenInt}$ inputs. Use
the native equality identifying $M$ with the $O$-span of the
projections of these generators.

The separating elements of LN3 give $V=\bigoplus_x V_x$.
Thus $\Phi(V)\subset\sum_x C_x$. On a rational form $y$, its
projection has $\Phi$-image equal to the component of $y$ along
$C_\emptyset$ in $\sum_x C_x$. This projection commutes with
each group-algebra action because both summands are stable. For
an integral conjugacy class sum $z$, apply the integral test to
each distinguished generator. Projecting the rational lattice
combination it gives shows that $z\Phi(1\otimes M)$ belongs to
$\Phi(1\otimes M)$: algebraic integers of the global coefficient
field map into $O$ at this place by the defining valuation.
Every central element of $O[G_{0g}]$ is an $O$-linear combination
of these integral class sums, since conjugation permutes the
group basis. Hence the full center preserves $\Phi(1\otimes M)$.
Pull its action back by the injective $m\mapsto\Phi(1\otimes m)$.
This defines an $O$-algebra action integrally, even if $\ell_i$
divides the group order. Rationalizing recovers the given action
on the image of $\Phi$ by semilinearity. It commutes with Hecke:
the central operator preserves each $C_x$, and consequently every
$V_x$ since $\Phi$ is injective and the different $C_x$ are
independent. Hecke acts by the specified point character on that
space. For any rational forms $a,b$ whose complex values are
related by the central action, project their equality along
$C_\emptyset$ to get the required equality of their projections
in $V$. If a rational form is in a single type space this same
projection lies in the indicated eigenspace or vanishes for an
absent point. Finally, if a central element acts by the chosen
$a\in O$ on every type space above $x$, it acts by its image on
all of $C_x$. Applying the injectivity of $\Phi$ gives the scalar
$a$ on the entire $V_x$. This is the original $\mathrm{EigSpan}$
quantifier. The construction costs no denominator from a group
order or a point projector.
% END M84T3

% BEGIN M84T4
\paragraph{M84T4. Rational and integral central operators in the selected row.}
Both the rational test $\mathrm{ZLStable}$ and the integral test
$\mathrm{CenInt}$ in M84T3 hold on the actual members. The
hypothesis of membership includes an $\mathrm{RXs}(t)$
witness. Its $\mathrm{AllT}$ realization identifies the finite
group $G$ with $G_{0g}$, its level-coordinate spaces with the
rational forms at $j=k+1$, and the rational form lattices with
the specified generator spans. For a central group-algebra
element over the global coefficient field, pull its coefficients
back by the group isomorphism. The resulting central operator
on level coordinates is defined over that field, and
$\mathrm{AllT}.\mathrm{cen}$ identifies its complex image with
the operator on forms. Applying it to the coordinates of any
rational form proves $\mathrm{ZLStable}$.

For the integral assertion let $E$ denote the inclusion matrix
from level coordinates into the ambient form space, and let
$L$ be the ambient integral lattice denoted by $\mathrm{nabLat}$.
The level lattice is $\{x:Ex\in L\}$. Right translations by the
finite group preserve $L$ in these row models: in the function
description the translated value at a class is obtained from
the value at its translated class by the weight action, which
carries the value lattice there onto the value lattice here.
This is the $\mathrm{nab}$ component of the integral action
test; the same change of class preserves the row Hermitian
pairing, which is positive definite at the chosen infinite
place. For $a=\sum_g n_g g$ with integral coefficients and $a$
central, $\sum_g n_g r(g)$ commutes with the idempotent imposing
the level character. Its image therefore preserves the level
subspace, and its level matrix $a_{\rm lev}$ satisfies
\[
       E a_{\rm lev}x=\sum_g n_g r(g)Ex\in L\qquad(Ex\in L).
\]
This proves integral preservation on level coordinates. Under
$\mathrm{AllT}.\mathrm{latL}$ each distinguished generator
belongs to this lattice (take its own coefficient to be one),
and every lattice output is a combination of the distinguished
complex generators over the global integers. Lift those
generators back by the rational coordinate isomorphism and
take the same combination. Its complex image equals the
central translate of the original generator by the operator
identity above. This is exactly $\mathrm{CenInt}$ in every
retained row.
% END M84T4

% BEGIN M84T5
\paragraph{M84T5. The two coefficient fields for the series idempotents.}
In the eventual quantifier defining membership take the fourth
finite list of prescribed algebraic complex numbers to contain
$\exp(2\pi i/N)$, where $N>0$ is the base's finite type-group
order. Its algebraicity follows from $X^N-1$. This list depends
only on the permitted base data, before $w,Z$. The global
coefficient field $L$ then contains a primitive $N$th root by
the defining inclusion test and injectivity of its complex
embedding. This is an allowed strengthening of the eventual
threshold. In the fraction field of every $O_i$ its image is
again primitive and is integral over $O_i$ by $X^N-1$. Since
$O_i$ is integrally closed, it lies in $O_i$ with the same
order. This supplies the exact $\mathrm{zetaPrimitiveRoot}$
entry for the desired original witness.

For either type prime $g_k$, $|\mathrm{GL}_r(\mathbb F_{g_k})|$
divides $N$ by the product formula for the type-group order:
the corresponding group order occurs to the positive power
$D.R.f$. In a complex irreducible representation $\tau$ of
that finite group every eigenvalue of $\tau(h^{-1})$ is an
$N$th root of unity. Thus its trace belongs to the complex
image of $L$. The coefficient at $h$ of the irreducible central
idempotent is
\[
      \frac{\dim\tau}{|\mathrm{GL}_r(\mathbb F_{g_k})|}
                          \operatorname{tr}\tau(h^{-1}).
\]
The member's $A.\mathrm{grp}$.series identity expresses the
specified idempotent as a sum of such idempotents. Its
coefficients therefore belong to the image of $L$. This proves
$\mathrm{SerLL}$ with the native coefficient field. Dividing
by the group order in this computation takes place in $L$.
% END M84T5

% BEGIN M84T6
\paragraph{M84T6. Five original Galois preconditions on the same member.}
The fields $\mathrm{large}$, $\mathrm{nct}$, $\mathrm{same}$,
$\mathrm{block}$ and $\mathrm{muUnr}$ of $\mathrm{GaloisPreconditions}(D,n)$ hold.
The coefficient condition is the $\mathrm{large}$ part of the
member's given nonempty $\mathrm{defaultExtension}$. The
norm-choice tie is the tie of the selected norm datum.

For $\mathrm{same}$, fix a localization $a$ over $w$. Its
actual $\mathrm{NdSpec}.\mathrm{tw\_reads}$ gives the
native homomorphism $f$ out of the initial coefficient ring
and embeddings of the two fraction fields into a field
$\Omega$. In $\Omega$ their readings agree on every element of
the initial ring and on every element of the global coefficient
field. The norm-choice tie gives a single global character
$\psi_0(w)$ whose evaluations at $g$ read as the two specified
local character values (with the prescribed Artin map on the
norm-choice side). Thus the images in $\Omega$ of every finite
sum $\sum_g f(a_g)\psi(g)$ are equal for the two readings.
The ring maps from the local rings into $\Omega$ are injective:
the small-to-large norm-choice ring map is injective, and all
remaining maps factor through inclusions into fraction fields
and field embeddings. Hence each such finite sum vanishes
at one reading exactly when it vanishes at the other. This
is $\mathrm{SameChar}$, and the index comparison in the
norm-choice tie gives $\mathrm{SameTwist}$ for each $a$.

For $\mathrm{block}$ put $j=D.\mathrm{nd}.j(i)-1$ as a
zero-based index and $e_\rho=(D.\mathrm{hc}\,\rho).e(j)$.
The member already carries the automorphic character tie
$c.a.\mathrm{chrN}$. It identifies
$\mathrm{tw}(\mu_C(i),\eta_T^{-1})$ with the tuple's
$j$th character times $\eta^{-1}$. The finite-order multiplier
of the tuple is trivial on the connected archimedean
components. Its infinity exponent is therefore that of the
underlying $\mu_0$, namely $-e_\rho$ in the chosen CM type.
The exponent of $\eta$ on that type is $1/2$. Inversion
negates exponents and multiplication adds them, as follows
directly by multiplying the character formulas on the
archimedean units. The two conjugate embeddings have opposite
angle exponents. The infinity-parameter clause of the same
character tie gives precisely the two singleton multisets
$\{-e_\rho-1/2\}$ and $\{e_\rho+1/2\}$ in the designated
order. These are the original $\mathrm{BlockCharParam}$ tests.

At a place $w$ above $\ell_i$ outside the two generator
primes, the same twisted rank-one automorphic parameter is
unramified. Use the unramified-place equivalence in
$c.a.\mathrm{chrN}$ to check the corresponding Hecke
character. The character $\eta$ is unramified here, and
$\mu_0$ can ramify only above a generator prime. In the tuple
multiplier $\delta\psi$, the chosen $\delta$ can ramify only
at its designated auxiliary prime. The norm choice uses the
other auxiliary character when $\ell_i$ equals the first
prime; the auxiliary primes are distinct. The character
$\psi$ is unramified outside its selected $Q$ at such a
place. The localization identifies this $Q$ with $Q_i$ by
$\mathrm{NdSpec}.Q_{\rm eq}$ and its base-$Q$ identity, and
the defining $S_i,Q_i$ condition excludes places above
$\ell_i$ from $Q_i$. Each factor is therefore trivial on
local units, as is their product. This proves
$\mathrm{MuUnr}$ with its coefficient-prime guard.
% END M84T6

% BEGIN M84T7
\paragraph{M84T7. Localized type sections from the member's Hecke realization.}
The section $\Phi$ used in M84T3 exists on the original member at
each $i$. Here are the Hecke inputs actually needed. At a split
place outside the excluded set, the member's $\mathrm{MtInput}$
gives hyperspecial level and trivial character there. Indeed
the two type-prime places are excluded, while the ramification
of the tuple multiplier is above the norm choice's $Q$ or the
residue characteristic, both excluded by the localization.
The chosen $\mathrm{NdSpec}$ has a rational operator on forms
realizing each of the specified spherical double cosets. This
also follows by taking $\mathrm{MtInput}.\mathrm{rat}$ in the
row $j-1$. Its projection is the specified local Hecke
operator. To see this, use the injective comparison $\Psi$
of the rationalized local lattice with the part of the
normalized Hecke space supported at its maximal ideal. On
the whole space the rational operator is the corresponding
generator, by its definition on rational forms. Its action
on the retained part is the localized action by
$\mathrm{NdSpec}.\mathrm{act\_hK}$. The projected rational
form differs from the original by a vector in the other
parts by $\Psi_{\rm pr}$. Those parts are stable under the
commuting Hecke generators. Thus the two candidate answers
differ by a vector in both the retained and the complementary
parts. These are disjoint by decomposition into the spaces
of distinct simultaneous characters. Injectivity of $\Psi$
gives the assertion.

The member's automorphic specialization gives precisely its
irreducible type spaces and their point test: equality of all
split-good Hecke scalars with those of $x$ is equivalent to
having point $x$. On each such type space each tested operator
is scalar. Indeed, at a split place the fixed space of a
hyperspecial subgroup in an irreducible spherical factor of
$\mathrm{GL}_{n_V}$ has dimension one, and the Hecke sum acts
only at this factor. All copies of the same representation
therefore have the same scalar. Equivalently, without choosing
a tensor factorization, the fixed space in the global
irreducible representation is a simple module for the
spherical Hecke algebra and the translations away from the
place. The translates of any nonzero vector generate the
irreducible representation, and averaging a finite sum of
such translates over the hyperspecial subgroup recovers the
fixed vectors using spherical double cosets. Spherical
convolution commutes (transpose preserves every double coset
by elementary divisors and reverses convolution), and it
commutes with the away translations. Its action on this
simple module is scalar: smoothness and second countability
bound the dimension of a cyclic representation by countable
dimension. A nonscalar endomorphism $T$ in its division
endomorphism ring would make $(T-a)^{-1}y$, $a\in\mathbb C$,
linearly independent for any $y\ne0$: multiply a finite
relation by the common denominator and evaluate the resulting
polynomial at each $a$. A nonzero polynomial in $T$ is
invertible since it splits into linear factors. This contradicts
countable dimension. Finally the level and hyperspecial Hecke
sums agree on these fixed vectors, since away from the place
one has the same support and at the place the same coset
representatives. The normalized scalar is the same on every
copy under a Hecke intertwiner.

Choose a basis of each $V_x$ in the direct sum of M84T3, with
coordinates $\lambda_{x,b}$. Each functional
$y\mapsto\iota_i\lambda_{x,b}(\mathrm{pr}\,y)$ on rational
forms extends to a complex-linear functional on $C$ by the
member's rational basis of forms, whose complex images are a
basis. Write $p:C\to\mathbb C^{\sum_x\dim V_x}$ for their
joint map. If the $x,b$ coordinate is nonzero on a vector
in a representative's type space, apply each of the tested
operators to that vector. On rational basis elements the
extended functional intertwines that operator with the
eigenvalue at $x$, by the projection just proved, hence it
intertwines on every vector. Cancelling the nonzero coordinate
identifies the representative's scalar with the scalar at $x$
for every split test. The point test then gives $x$ as its
point. Hence $p$ takes $C_x$ into its own coordinate block
and kills $C_\emptyset$.

The projections of rational forms span $V$ over $K$.
Consequently the image of $p$ contains the coordinate
basis, and $p$ is surjective. Since $C$ is spanned by the
type spaces, its restriction to $G=\sum_x C_x$ is still
surjective. The member's multiplicity identity gives
$\dim C_x=\dim_K V_x$, so $\dim G\le\sum_x\dim_K V_x$.
Thus $p|_G$ is an isomorphism. Take its inverse on the
embedded $K$-coordinates of $V$. The resulting map is
injective, semilinear, and takes the $x$-space into $C_x$
(the block images already span the individual blocks).
Its composition with projection on rational forms differs
from the original complex image by $\ker p=C_\emptyset$.
These are exactly all fields of $\mathrm{Sec}$.

In particular the remaining hypothesis in M84T3 is the
preservation of the individual type spaces by the action.
It suffices to identify the complex central operator on the
actual realization with its sum of adelic right translates
(the member comparison $\mathrm{CenLoc}$): each isotypic
sum in this realization is right-translation stable. On
noncentral elements the native $\mathrm{onMt}$ is zero.
% END M84T7

% BEGIN M84T8
\paragraph{M84T8. From local parahoric vectors to the fourth original tie.}
Fix a series system $S_k$ for each type prime whose rank-$r$
projectors are the original $e_{\rm Ser}(f)$. Assume the
local representation comparison $\mathrm{LocType}(D,S)$
at both primes. It states that the nonzero central actions
on type vectors have parahoric block vectors in a single
local representation $\pi$, supplies the determinant-character
label on the lower block and the matching inertia value of
$\mu_C$, and identifies the inertia polynomial of the
automorphic parameter from any full-rank series detected in
$\pi$. The member satisfies the additional unramified-twist
identity $\mathrm{IneUnr}(D)$: its $A.\mathrm{wdT}$ carries
rank-one twist to tensor product and preserves unramifiedness,
dimension and inertia polynomials. On every inertia element
the unramified rank-one factor acts identically, so tensoring
it with any matrix preserves the characteristic polynomial.
Thus the full $\mathrm{SerLoc}(D)$ follows from the stated
series system and local representation comparison.

Fix a parameter $\sigma$ for which $e_{\rm Ser}(f)$ acts
nontrivially on a vector in its parameter space. The native
span identity for the parameter space then gives a single
type space on which this action is nonzero: otherwise its
linear kernel would contain their span. Choose its actual
representative $P$ and the list $\Sigma$ in the span test.
Its isobaric localizations give $\sigma$, and appending the
distinguished rank-one string and twisting by $\eta_T^{-1}$
is a base change of $P$. The local comparison gives $\pi$,
$\chi,c$, and a parahoric vector with nonzero top projection.
In the principal-congruence fixed space write $u$ for that
projected vector. It is fixed by the parabolic unipotent
radical: translating by a top-block matrix conjugates the
radical into itself. The top projector fixes $u$ by
idempotence. The lower-block projector for
$h=(X-c)^t$ acts on the original vector by a nonzero
idempotent scalar on the determinant character $\chi$,
hence by one; it commutes with the top action, so fixes
$u$ as well. The product property in the series system
therefore gives $e_{n_V}(fh)u=u\ne0$.

Write $[f]_\tau$ for the product of $X-\tau(x)$ over
the roots of $f$ in an algebraic closure of the residue
field, with multiplicities, where $\tau$ is the specified
injective Teichm\"uller map. The full-rank comparison gives
$[fh]_\tau$ for the inertia polynomial on $P$ at the
specified inertia generator. On the base-change list this
same polynomial is
\[
       [\mathrm{labPar}(\sigma)]_\tau\,(X-\tau(c))^t.
\]
Indeed sums multiply inertia polynomials, the empty sum has
polynomial one, and the real unramified shifts in a Speh
string have trivial inertia, so a string of length $d$
contributes the $d$th power. The rank-one $\eta_T^{-1}$
twist disappears by the unramified-twist test. The
degree of the parameter's label equals its rank, so the
native degree guard in reading its inertia polynomial
holds. The lower-block value is precisely the scalar
prescribed by $\mathrm{LocType}$ for $\mu_C$. Cancel
$(X-\tau(c))^t$ in $\mathbb C[X]$. The remaining
root multisets agree by injectivity of $\tau$, including
all possible repeated roots. The two polynomials are monic
(the first by the nonzero series projector and the second by
the parameter condition); factorization over the residue
algebraic closure and injectivity of base change now give
$f=\mathrm{labPar}(\sigma)$. This proves the asserted
series label for every vector in the original test. The
comparison $\mathrm{LocType}$ itself concerns irreducible
local representations and depth-zero types; M84T8 isolates
it as an input.
% END M84T8

% BEGIN M84T9
\paragraph{M84T9. The same tame inertia characters at all original type places.}
On the original member one can choose the entire
$\mathrm{om}$ family of the desired $\mathrm{Pend}$ witness,
including its eigenvalue and character-twist equalities.
The initial $A.\mathrm{lab}$ reads the Weil--Deligne label
at each selected generator $\gamma_{k,v}$. The given
$\mathrm{defaultExtension}$ supplies its conjugate-place
counterpart $\mathrm{labc}$, including a lift of CM conjugation
identifying the valuation rings and conjugating the generators.
At a selected place take the given generator, and at its
conjugate take the generator in this latter identification.
These assignments agree on their domains: above a real place
of the completely split type prime there are exactly two
places, exactly one in the chosen half. The selected-place
index is injective. Complete the assignment by the identity
away from the type prime. This proves the member's
$\mathrm{LabAll}$ and transports tame generation to every
conjugate place: the indicated conjugacy carries inertia,
its wild subgroup and all their open subgroups to the
corresponding groups on the other side.

For such a place $w$ over $g=g_k$, the residue field has
order $g$. The tame inertia quotient is procyclic of
prime-to-$g$ order, generated topologically by the chosen
$\gamma_w$. More explicitly over the maximal unramified
extension, adjoining $d$th roots of a uniformizer gives
the cyclic quotients of degree $d$ prime to $g$: all
$d$th roots of unity are already present by residue-field
lifting, and adjoining the root has degree $d$ by valuation.
These quotients exhaust the tame inertia action. For
$x\in\overline{\mathbb F}_g^\times$ take its finite order
$d$, which is prime to $g$. There is a unique tame character
of inertia with value $\tau(x)$ at $\gamma_w$: use the cyclic
quotient of order $d$, and uniqueness follows from tame
generation. Denote it by $\omega_{w,x}$, with $\tau$ the
specified multiplicative lift. Set the value to one off
inertia; use the trivial character for $x=0$ and at places
away from $g$. These assignments satisfy all the native
one, multiplicativity, wild, and openness tests. In detail
multiplication in $x$ follows by uniqueness on the generator,
and finite image gives an open kernel in inertia. The
induced profinite topology provides an open subgroup of the
global Galois group whose intersection with inertia lies in
that kernel. Distinct nonzero $x$ are distinguished already at
$\gamma_w$ by injectivity of $\tau$.

For any local parameter, the label records the eigenvalues
of this generator on its wild-invariant subspace, with
multiplicity. Indeed inertia on a complex Weil--Deligne
representation has finite image, and on the wild-invariant
part the generator has prime-to-$g$ order. Its eigenvalues
therefore have unique preimages under $\tau$; the multiset
is stable under $x\mapsto x^g$ by Frobenius conjugation
on tame inertia. It gives the unique monic polynomial over
$\mathbb F_g$ with that multiset of roots, exactly the
Weil--Deligne label in $\mathrm{LabAll}$. Its degree is
the dimension of this wild-invariant subspace. Under the
native full-degree guard the whole space is wild-invariant.
The finite tame action is a direct sum of eigenspaces of
the cyclic quotient. Its characters are uniquely determined
by the generator, hence are the $\omega_{w,x}$ just chosen
with the multiplicities of the label's roots. This proves
the entire polynomial identity for every inertia element.
It also covers the zero-dimensional representation, for
which the empty product is one.

Finally, suppose a rank-one parameter has label $X-u$ with
$u\in\mathbb F_g^\times$. It too has full degree, and its
inertia character is $\omega_{w,u}$. The member's
$A.\mathrm{wdT}$ identifies the rank-one twist with tensor
product on Weil--Deligne representations. Tensoring the
previous tame representation by this line stays tame of the
same dimension. Each root $x$ of the original label is
multiplied by $u$, because $\tau(ux)=\tau(u)\tau(x)$ and
the tame characters agree on their generator. This proves
the degree and multiset assertions for the twisted parameter.
Thus the very same choice supplies both
$\mathrm{inertiaCharpoly\_eigenvalues}$ and
$\mathrm{inertiaCharpoly\_character}$ on all original places.
The remaining local comparison in M84T8 and the rank-one
de Rham and depth-zero representation tests in the first
two clauses of $\mathrm{Tie4}$ have separate content.
% END M84T9

% BEGIN M84T10
\paragraph{M84T10. The first character tie on the original member.}
Fix an original coefficient index $i$, a generator prime $p=g_k=\ell_i$
and $v\in\mathrm{PlG}(k)$. Write $w_+=\mathrm{plGw}(k,v)$,
$w_-=\overline{w_+}$ and $j=j(i)-1$ in zero-based indexing.
Both completions are $\mathbb Q_p$ by complete splitting. We prove
$\mathrm{tieChi}$ using global and local Artin reciprocity, Chebotarev
density and the rank-one crystalline period calculation.

Let $\mu=\mathrm{tp.tup}(\mu_i,j)$, $\lambda=\mu\eta^{-1}$
and $s=(t-1)/2+(1-n_V)/2=-r/2$. At every split good prime $z$
outside the finite ramification set and away from $p$, the member's
$\mathrm{chrN}$ identifies the one Hecke eigenvalue of this line
with $\lambda(\varpi_z)^{-1}$. The Frobenius condition on
$\chi_1=\mathrm{nd.chi}(i,1)$ therefore reads
\[
 \iota_i\bigl(\chi_1(\mathrm{Frob}_z)\bigr)
        =\lambda(\varpi_z)^{-1}(\mathrm Nz)^s .
\]
Here the embedding on the left includes the indicated embedding of the
coefficient ring. Put $\beta=\lambda|\cdot|^s$, using the idelic
absolute value, so the right side is $\beta_z(\varpi_z)^{-1}$.
The angle exponents of M84T6 give the two infinity exponents of
$\beta$ above each real embedding as
$s-(e_\rho+1/2)$ and $s+(e_\rho+1/2)$. Both are integers because
$e_\rho\in\mathbb Z$ and $r$ is odd. Thus $\beta$ is an algebraic
Hecke character.

For precision take the inverse finite character $\beta_f^{-1}$ and
the Artin normalization sending uniformizers to the Frobenius in this
display. Its $p$-adic avatar can be read directly on finite ideles.
On a principal idele $x$, $\beta_f^{-1}(x)=\beta_\infty(x)$ is
the product of the indicated integral embedding powers. Hence all
values of $\beta_f^{-1}$ are algebraic and lie in a fixed number
field: the finite local unit characters have finite image, and
taking powers of finitely many ideal class representatives reduces
their values to principal values and unit values. Transport
$\beta_f^{-1}$ through $\iota_i^{-1}$ and divide by the product
of these same algebraic powers on the $p$-components, grouping the
embeddings by the places they induce. The result is trivial on
$E^\times$. It is continuous: at $p$ the correction is locally
algebraic and the finite part is locally constant, while away
from $p$ only finitely many unit groups act nontrivially.
It therefore factors through the profinite idele class quotient
by Artin reciprocity. Indeed modulo the closure of $E^\times$ in
the finite ideles the unit groups have compact image of finite
index, by finiteness of the ideal class group; the resulting
quotient is profinite. At a good $z\nmid p$ it takes the displayed
Frobenius value. These Frobenius elements determine a continuous
abelian character: given a finite Galois quotient over $E$, take
a Galois closure over $E^+$ and lift the desired element to the
subgroup fixing $E$. Chebotarev over $E^+$ gives a prime outside
the finite exclusions, splitting in $E$, and a place above it with
that Frobenius in the quotient. This proves equality of the avatar
with the original $\chi_1$.

At either $w=w_\pm$, the correction on local units is an integral
power of the identity under $E_w=\mathbb Q_p$. Local reciprocity
therefore expresses $\chi_1$ as a Tate power times a finite
character and an unramified character. Explicitly extend the smooth
unit character to the local multiplicative group by assigning one
to $p$, and remove this finite character and the indicated
cyclotomic power from $\chi_1$; the quotient is unramified.
The unramified line is crystalline by the rank-one period
calculation of S80J4; the Tate line has a period given by a power
of Fontaine's $t$. Their Weil--Deligne inertia is trivial.
After a finite extension killing the smooth character the remaining
finite line has the period $1\otimes e$. Inertia acts on this
potentially crystalline period through the original finite
character, since inertia fixes unramified period constants.
Finite-extension descent gives de Rhamness on $E_w$: after normal
closure the full de Rham periods descend by semilinear Galois
descent. The member's $\mathrm{loc\_tie}$ and M84T1 read this on
the required $\mathrm{nbr}$ and identify the native inertia
polynomial. It remains to identify the smooth character on one
specified generator at each place.

On units of $w_+$ the factors $|\cdot|^s$ and $\eta$ are trivial.
In $\mu=\mu^0_j\xi$ the multiplier is unramified there: its
selected $\delta$ can ramify only at its selected auxiliary prime,
distinct from the coefficient prime by the choice in M84T6, and
$\psi$ is trivial above that coefficient prime. This includes the
reference tuple with trivial multiplier. Thus the smooth unit
character is read from the tame character $\mu^0_j$, with the
fixed inverse from the avatar. Let $\gamma$ be the designated
tame generator and choose $\varpi\in w_+\setminus w_+^2$ in
$\mathcal O_E$. For $\pi^{p-1}=\varpi$, the residue of
$\gamma(\pi)/\pi$ belongs to $\mathbb F_p^\times$. Lift it to
$b\in\mathcal O_E$, necessarily a unit at $w_+$. This is exactly
the exponent-one test $\mathrm{TameRep}(w_+,1,\gamma,b)$.
The local Artin unit of $\gamma$ has residue $\bar b$ or
$\bar b^{-1}$ according to reciprocity convention. Indeed
$\pi/(\zeta_p-1)$ is an algebraic unit, whose residue inertia
fixes, and
$\gamma(\zeta_p-1)/(\zeta_p-1)$ reduces to the mod-$p$
cyclotomic character of $\gamma$. Local cyclotomic reciprocity on
$\mathbb Q_p^\times$ sends a unit to that unit or its inverse
on $p$-power roots of unity. The member's $A.\mathrm{chiT}$
now gives $\mu^0_{j,w_+}(b)=\tau(c)$, with
$c\in\{\chi_J,\chi_J^{-1}\}$. Since the character is tame, only
the residue of the Artin unit matters. Its value on $\gamma$ in
the Weil--Deligne line is consequently $\tau(u)$ for one
$u\in\{\chi_J,\chi_J^{-1}\}$, absorbing the fixed inversions.

Finally $\mu^c=\mu^{-1}$ on conjugate local units: apply the
$\mathrm{csd}$ equivalence in $\mathrm{chrN}$ to
$\mathrm{muC\_csd}$ and use its tuple equality, on finite ideles supported at the tested place. For the local
component first take $x\in E^\times$ as in the definition of
$\mathrm{csd}$, then extend to $E_w^\times$ by density.
M84T9 chooses the conjugate generator by transport, and local
reciprocity commutes with this transport. Thus the smooth
Weil--Deligne character at $w_-$ takes the value $\tau(u^{-1})$
there. Both inertia actions are tame on the full line; their native
labels are therefore $X-u$ and $X-u^{-1}$ by M84T9. These are
distinct places; for every other $w\mid p$ the two implications
in the definition are vacuous. This proves the first original
tie, with the same $u$ for both tests.
% END M84T10
